\documentclass[11pt]{article}
\usepackage[a4paper,margin=1in]{geometry}
\usepackage{amsmath,amsthm,amssymb}
\usepackage{mathtools}

\newtheorem{theorem}{Theorem}[section]
\newtheorem{proposition}[theorem]{Proposition}
\newtheorem{lemma}[theorem]{Lemma}
\newtheorem{corollary}[theorem]{Corollary}
\theoremstyle{definition}
\newtheorem{definition}[theorem]{Definition}
\newtheorem{remark}[theorem]{Remark}

\newcommand{\R}{\mathbb{R}}
\newcommand{\Fr}{\mathcal{A}}
\newcommand{\Om}{\Omega}
\newcommand{\eps}{\varepsilon}
\newcommand{\rstar}{\rho_{*}}
\newcommand{\zbar}{\bar z_{\rho}}
\newcommand{\er}{e_{r,\zbar}}

\title{H$^{1}$-in-$\rho$ Bound for the Centroid of the Torsion Level Family\\
       (Route~$\delta$, Plan~2, Building Block \textsc{CentroidBound})}
\author{}
\date{}

\begin{document}
\maketitle

\begin{abstract}
This note is the building block \textsc{CentroidBound} of the Plan~2
proof of sharp quantitative Faber--Krahn stability via Route~$\delta$
(no graph, no selection, no Schauder). For a near-spherical torsion
domain $\Om\subset B_R$ with $|\Om|=\omega_n$ and small Saint--Venant
deficit $\delta:=\delta_T(\Om)$, let $E_\rho:=\{u>t(\rho)\}$ be the
volume-radius parametrised superlevels of the torsion function $u$
and let
\[
   \zbar:=\frac{1}{|E_\rho|}\int_{E_\rho}x\,dx\in\R^n
\]
be the \emph{centroid} of $E_\rho$. We prove that $\rho\mapsto\zbar$
is Lipschitz on $[\rstar,1]$ and satisfies the integrated
$H^{1}$-in-$\rho$ bound
\[
   \int_{\rstar}^{\rho_\delta}|\bar z_\rho'|^{2}\,(-t'(\rho))\,d\rho
   \;\le\; C(n,R,\rstar)\,\delta_T(\Om).
\]
The proof uses the centroid kinematic identity (a bulk-coarea
computation that has \emph{no} homothety term), the unconditional
Cauchy--Schwarz of the velocity defect controlled by $D_H+D_I$, the
Figalli--Maggi--Pratelli volume Fraenkel estimate, the Fusco--Julin
$L^1$ oscillation, the divergence theorem for finite-perimeter sets,
and the integrated Talenti gap of \cite{LevelSetIdentity}. No graph parametrisation, no
selection principle, no pointwise quantitative isoperimetric form is
invoked.
\end{abstract}

\section{Statement and main result}\label{sec:statement}

Fix $n\ge 2$, $R\ge 2$, $\rstar\in(0,1)$, and $\kappa>0$. Let
$\Om\subset B_R\subset\R^n$ be a bounded open set with
$|\Om|=\omega_n:=|B_1|$ and Saint--Venant deficit
$\delta:=\delta_T(\Om)\le\delta_0(n,R,\rstar)$, where
$\delta_T(\Om):=1-T(\Om)/T(B_1)$ and $T(\cdot)$ denotes the
Saint--Venant torsional rigidity. Let $u\in
W^{1,2}_0(\Om)\cap W^{2,p}_{\rm loc}(\Om)$ ($p>n$) be the variational
torsion function of $\Om$, $t:[\rstar,1]\to(0,\|u\|_\infty)$ the
volume-radius parametrisation defined by
$|\{u>t(\rho)\}|=\omega_n\rho^n$, $E_\rho:=\{u>t(\rho)\}$, and
\[
   V_\rho(x):=\frac{-t'(\rho)}{|\nabla u(x)|}
   \qquad\text{on }\partial^*E_\rho\cap\{|\nabla u|>0\}.
\]
The set $\{|\nabla u|=0\}\cap\partial^*E_\rho$ has
$\mathcal H^{n-1}$-measure zero on $G_{\rm CL}$ (Sard plus regularity,
cf.\ \cite{LevelSetIdentity}), so $V_\rho$ is
$\mathcal H^{n-1}$-a.e.\ defined on $\partial^*E_\rho$ at good levels.
Set $\rho_\delta:=1-\kappa\sqrt{\delta}$.

\begin{definition}[Centroid field]\label{def:centroid}
For $\rho\in[\rstar,1]$ define the \emph{centroid} of $E_\rho$ by
\begin{equation}\label{eq:centroid-def}
   \zbar:=\frac{1}{|E_\rho|}\int_{E_\rho}x\,dx\in\R^n.
\end{equation}
This is the bulk barycentre; it is \emph{not} obtained by Fraenkel
minimisation and inherits no implicit-function-theorem regularity.
\end{definition}

\begin{theorem}[$H^{1}$-in-$\rho$ centroid bound]\label{thm:F}
There exist $\delta_0=\delta_0(n,R,\rstar)>0$ and
$C_{\rm cent}=C_{\rm cent}(n,R,\rstar)>0$ such that for every bounded
open $\Om\subset B_R$ with $|\Om|=\omega_n$ and $\delta\le\delta_0$,
the centroid field $\rho\mapsto\zbar$ is Lipschitz on $[\rstar,1]$
with explicit constant
\begin{equation}\label{eq:Lcent}
   L_{\rm cent}\;:=\;\frac{2nR}{\rstar^{n}},
\end{equation}
in particular absolutely continuous, and
\begin{equation}\label{eq:F}
   \int_{\rstar}^{\rho_\delta}
     |\bar z_\rho'|^{2}\,(-t'(\rho))\,d\rho
   \;\le\;C_{\rm cent}\,\delta_T(\Om). \tag{F}
\end{equation}
\end{theorem}

\begin{remark}[Compatibility with the Fraenkel centre]\label{rmk:CFr}
By Cicalese--Leonardi \cite[Prop.~2.3]{CL2012} and Acerbi--Fusco--Morini
\cite[Lem.~2.6]{AFM2013}, in the smallness regime
$|\zbar-z^{\rm Fr}_\rho|\le C\,\mathcal A(E_\rho)\le C'\sqrt{\delta_T(\Om)}$.
The centroid is therefore interchangeable with $z^{\rm Fr}_\rho$ in
downstream uses at cost $O(\sqrt{\delta_T})$, absorbed into
constants exactly as in Wave~2~C~§C1.
\end{remark}

\section{Setup, measurability, Lipschitz-in-$\rho$}\label{sec:setup}

\subsection{Good levels and the CL regime}

The exact level-set identity gives the integrated bound
\begin{equation}\label{eq:int-DHDI}
   \int_{\rstar}^{1}\bigl(D_H(t(\rho))+D_I(t(\rho))\bigr)\,d\mu(\rho)
   \le C(n,\rstar)\,\delta_T(\Om),
   \qquad d\mu=(-t')\,d\rho.
\end{equation}
It does \emph{not} imply pointwise smallness of \(D_I(t(\rho))\) for
every \(\rho\). We therefore use the following good-level split. Let
\(\eta_{\rm CL}(n,\rstar)>0\) be small enough that
\[
   D_I(t(\rho))\le\eta_{\rm CL}
   \quad\Longrightarrow\quad
   \mathcal A(E_\rho)\le\eps_0(n)
\]
by the quantitative isoperimetric inequality. Define
\[
   G_{\rm CL}:=\{\rho\in[\rstar,1]:D_I(t(\rho))\le\eta_{\rm CL}\},
   \qquad B_{\rm CL}:=[\rstar,1]\setminus G_{\rm CL}.
\]
Then \(\mu(B_{\rm CL})\le C\delta_T(\Om)\) by
\eqref{eq:int-DHDI}. (Chebyshev: on $B_{\rm CL}$ one has $D_I(t(\rho))
>\eta_{\rm CL}$, so $\eta_{\rm CL}\,\mu(B_{\rm CL}) \le
\int_{B_{\rm CL}}D_I\,d\mu \le \int_{\rstar}^{1}(D_H+D_I)\,d\mu
\le C(n,\rstar)\,\delta_T(\Om)$, whence
$\mu(B_{\rm CL})\le C'\delta_T(\Om)$.)
On \(G_{\rm CL}\), Cicalese--Leonardi gives a
Fraenkel centre \(z^{\rm Fr}_\rho\) with
\begin{equation}\label{eq:CL}
   B(z^{\rm Fr}_\rho,\rho/2)\subset E_\rho\subset B(z^{\rm Fr}_\rho,2\rho).
\end{equation}
Combined with the centroid--Fraenkel proximity, after decreasing
\(\delta_0\) if needed,
\[
   B(\zbar,\rho/4)\subset E_\rho\subset B(\zbar,3\rho)
   \qquad\text{for a.e. }\rho\in G_{\rm CL}.
\]
In particular \(|x-\zbar|\le 3\) on \(\partial^*E_\rho\) and
\(\zbar\in E_\rho^\circ\) on the good levels. On \(B_{\rm CL}\) we use
only the global Lipschitz bound of Lemma~\ref{lem:Lip}.

\subsection{Borel measurability and Lipschitz}

\begin{lemma}\label{lem:Lip}
$\rho\mapsto\zbar$ is Borel and Lipschitz on $[\rstar,1]$ with
constant $L_{\rm cent}=2nR/\rstar^{n}$.
\end{lemma}

\begin{proof}
\emph{Borel.} Since $t$ is decreasing, $\rho\mapsto E_\rho$ is monotone
nested; dominated convergence gives $\rho\mapsto\int_{E_\rho}x_i\,dx$
continuous, hence Borel.

\emph{Lipschitz.} For $\rstar\le\rho\le\rho'\le 1$, nesting gives
\[
   |E_\rho\Delta E_{\rho'}|=|E_{\rho'}|-|E_\rho|
   =\omega_n((\rho')^n-\rho^n)\le n\omega_n|\rho-\rho'|.
\]
Since $E_\rho,E_{\rho'}\subset B_R$,
$|\int x_i-\int x_i|\le Rn\omega_n|\rho-\rho'|$. Combined with
$|E_\rho|\ge\omega_n\rstar^n$,
$|\zbar-\bar z_{\rho'}|\le (2nR/\rstar^n)|\rho-\rho'|$.
\end{proof}

\section{The centroid kinematic identity}\label{sec:kinematic}

\begin{lemma}[Centroid kinematic formula]\label{lem:3.1}
For a.e.\ $\rho\in[\rstar,1]$,
\begin{equation}\label{eq:3.2}
   \bar z_\rho'
   =\frac{1}{|E_\rho|}\int_{\partial^*E_\rho}(x-\zbar)\,V_\rho\,d\mathcal H^{n-1}.
\end{equation}
\end{lemma}

\begin{proof}
By the coarea identity \cite[\S3]{LevelSetIdentity} and Sobolev--Sard
\cite[App.~A, Thm.~A.1]{FigalliMaggi2011}, for every Lipschitz $f:B_R\to\R$ and
a.e.\ $\rho\in[\rstar,1]$,
$\frac{d}{d\rho}\int_{E_\rho}f\,dx=\int_{\partial^*E_\rho}f V_\rho\,d\mathcal H^{n-1}$.
By the volume normalisation $|E_\rho|=\omega_n\rho^n$ we have
$\tfrac{d}{d\rho}|E_\rho|=n\omega_n\rho^{n-1}=P(B_\rho)$, and combining
with coarea ($f\equiv 1$) gives
$\int_{\partial^*E_\rho}V_\rho\,d\mathcal H^{n-1}=\tfrac{d}{d\rho}|E_\rho|=P(B_\rho)$.
Taking $f(x)=x_i$ in the coarea identity and applying the product
rule to $\bar z_{\rho,i}=|E_\rho|^{-1}\int_{E_\rho}x_i\,dx$:
\[
   \bar z_{\rho,i}'
   =\frac{1}{|E_\rho|}\int_{\partial^*E_\rho}x_i V_\rho\,d\mathcal H^{n-1}
   -\frac{1}{|E_\rho|^2}\Bigl(\tfrac{d}{d\rho}|E_\rho|\Bigr)
       \int_{E_\rho}x_i\,dx.
\]
Substituting the coarea form
$\tfrac{d}{d\rho}|E_\rho|=\int_{\partial^*E_\rho}V_\rho\,d\mathcal H^{n-1}$
into the second term yields
\[
   \bar z_{\rho,i}'
   =\frac{1}{|E_\rho|}\int_{\partial^*E_\rho}(x_i-\bar z_{\rho,i})\,V_\rho\,d\mathcal H^{n-1},
\]
the displayed identity. The equality
$\tfrac{d}{d\rho}|E_\rho|=P(B_\rho)$ is here a \emph{consequence} of
the volume normalisation, not a perimeter identity for $E_\rho$.
\end{proof}

\begin{remark}[No homothety term --- the structural mechanism]
\label{rmk:no-homothety}
Identity~\eqref{eq:3.2} is the first vector moment of $V_\rho$ relative
to the centroid. \emph{No homothetic vector field
$H_{z,\rho}(x)=(x-z)\cdot\nu_E/\rho$ enters.} By contrast, defining
the centre field via the Fraenkel functional first-order condition
would yield a kinematic formula whose right-hand side involves
$V_\rho-H_{z,\rho}$, and the bound would consume the per-$\rho$
homothetic velocity defect (audit-trail markdown note, input~(2.4)),
itself conditional on the per-$\rho$ radial trace~(R)/(1.2), which is
open in the no-graph regime. The centroid choice eliminates the
homothety because it is defined directly by a bulk integral, not by a
minimisation principle comparing to balls. This is precisely the
structural mechanism by which Theorem~\ref{thm:F} avoids (R), (B²),
and the per-$\rho$ pointwise radial-trace inputs of the Wave-1 audit
(audit-trail markdown note, inputs~(1.1), (1.2)=(R), (2.4) in the
source notation) entirely.
\end{remark}

\section{Pointwise bound $|\bar z_\rho'|^{2}\le C(D_H+D_I)$}\label{sec:pointwise}

We bound the integrand in~\eqref{eq:F} pointwise by the two algebraic
deficits $D_H(t(\rho)),D_I(t(\rho))$ of \cite{LevelSetIdentity}.

\begin{lemma}[Pointwise centroid-velocity bound on good levels]\label{lem:4.1}
There is $C_*=C_*(n,R,\rstar)$ such that for a.e.\ $\rho\in G_{\rm CL}$,
\begin{equation}\label{eq:4.1}
   |\bar z_\rho'|^{2}\le C_*\bigl(D_H(t(\rho))+D_I(t(\rho))+\delta_T(\Om)\bigr).
\end{equation}
\end{lemma}

\begin{proof}
Set $\bar V_\rho:=P(B_\rho)/P(E_\rho)\in(0,1]$. Decompose
$V_\rho=\bar V_\rho+(V_\rho-\bar V_\rho)$ in~\eqref{eq:3.2}:
\begin{equation}\label{eq:4.2}
   |E_\rho|\bar z_\rho'
   =\underbrace{\bar V_\rho\!\int(x-\zbar)d\mathcal H^{n-1}}_{(\mathrm I)}
   +\underbrace{\int(x-\zbar)(V_\rho-\bar V_\rho)d\mathcal H^{n-1}}_{(\mathrm{II})}.
\end{equation}

\emph{Bound on (II).} By CL, $|x-\zbar|\le 3$, so
$|(\mathrm{II})|\le 3\int|V_\rho-\bar V_\rho|\,d\mathcal H^{n-1}$.
Set $A_\rho:=\int_{\partial^*E_\rho}V_\rho\,d\mathcal H^{n-1}=P(B_\rho)$
(by coarea with $f\equiv 1$ as in the proof of Lemma~\ref{lem:3.1});
the mean here is the perimeter-mean
$\bar V_\rho=A_\rho/P(E_\rho)$. The algebraic Cauchy--Schwarz
\[
   \Bigl(\int_{\partial^*E_\rho}|V_\rho-\bar V_\rho|\,d\mathcal H^{n-1}\Bigr)^2
   \;\le\; P(E_\rho)\,
     \int_{\partial^*E_\rho}(V_\rho-\bar V_\rho)^2\,d\mathcal H^{n-1}
\]
together with the variance-form bound
$\int(V_\rho-\bar V_\rho)^2\,d\mathcal H^{n-1}\le C(n,\rstar)\,D_H(t(\rho))$
(exact identity of \cite{LevelSetIdentity}, controlling $D_H$ by the
$L^2$ variance of $V_\rho$ around $\bar V_\rho$), and the perimeter
bound $P(E_\rho)\le C(n,\rstar)$ on $G_{\rm CL}$, give
\begin{equation}\label{eq:4.3}
   \Bigl(\int|V_\rho-\bar V_\rho|\,d\mathcal H^{n-1}\Bigr)^2
   \le C_{n,\rstar}\,D_H(t(\rho)).
\end{equation}

\emph{Bound on (I).} We claim, for each $i$,
\begin{equation}\label{eq:4.4}
   \Bigl|\int_{\partial^*E_\rho}(x-\zbar)_i\,d\mathcal H^{n-1}\Bigr|
   \le C_2(n,\rstar)\bigl(\sqrt{\delta_\rho}+\sqrt{\delta_T}\bigr),
\end{equation}
with $\delta_\rho:=P(E_\rho)-P(B_\rho)$. To see this, write
$(x-\zbar)_i=|x-\zbar|(e_r)_i$, and split
$(e_r)_i=(\nu_E)_i+((e_r)_i-(\nu_E)_i)$.

For the first piece, apply the divergence theorem to the Lipschitz
field $V^{(i)}(x):=|x-\zbar|\,e_i$ (where $e_i$ is the $i$-th standard
basis vector, so $V^{(i)}$ has only its $i$-th component nonzero).
Hence
\[
   \mathrm{div}\,V^{(i)}(x)
   = \partial_i\,|x-\zbar|
   = (e_r(x))_i,
   \qquad e_r(x):=\frac{x-\zbar}{|x-\zbar|}.
\]
The kink at $x=\zbar$ is a single point; since $n\ge 2$ this set has
codimension $\ge 2$, hence is $\mathcal H^{n-1}$-null and contributes no
singular part to the distributional gradient. By
\cite[Thm.~15.9]{Maggi2012},
\begin{equation}\label{eq:4.6}
   \int_{\partial^*E_\rho}|x-\zbar|(\nu_E)_i d\mathcal H^{n-1}
   =\int_{E_\rho}(e_r(x))_i\,dx.
\end{equation}
On $B_\rho(\zbar)$, $\int(e_r)_i=0$ by spherical symmetry. Hence
$|\int_{E_\rho}(e_r)_i|\le|E_\rho\Delta B_\rho(\zbar)|$. By FMP
\cite{FMP2010} at the Fraenkel centre $z^{\rm Fr}_\rho$,
$|E_\rho\Delta B_\rho(z^{\rm Fr}_\rho)|\le C\sqrt{\delta_\rho}$, while
the centroid--Fraenkel proximity
$|\zbar-z^{\rm Fr}_\rho|\le C'\sqrt{\delta_T}$ (Remark~\ref{rmk:CFr})
costs an additional
$|B_\rho(\zbar)\Delta B_\rho(z^{\rm Fr}_\rho)|
\le |\zbar-z^{\rm Fr}_\rho|\cdot P(B_\rho)
\le C''\sqrt{\delta_T}$
in recentring. Hence
\[
   |E_\rho\Delta B_\rho(\zbar)|
   \le C\bigl(\sqrt{\delta_\rho}+\sqrt{\delta_T}\bigr).
\]
On the good set $G_{\rm CL}$, both $\delta_\rho$ and the recentring
cost $\sqrt{\delta_T}$ are controlled, so the rate is preserved.

For the second piece, $|x-\zbar|\le 3$ and the FJ $L^1$ oscillation
\cite{FJ2014} at the centroid (with $O(\sqrt{\delta_\rho})$ change-of-centre
error absorbed): $\int|\nu_E-e_r|\le C_{\rm FJ}\sqrt{\delta_\rho}$.

Combining gives \eqref{eq:4.4}.

\emph{Combining.} From \eqref{eq:4.2}, \eqref{eq:4.3}, \eqref{eq:4.4}:
\[
   |E_\rho|^2|\bar z_\rho'|^2
   \le 4C_2^2\,\delta_\rho + 4C_2^2\,\delta_T + 2C_1 D_H,
\]
where $C_1=9C_{n,\rstar}$ (the $9$ comes from $3^2$ since
$|x-\zbar|\le 3$ on the bounded reduction). Use
$\delta_\rho\le D_I/(n\omega_n\rstar^{n-1})$ and
$|E_\rho|^2\ge\omega_n^2\rstar^{2n}$ to get \eqref{eq:4.1}.
\end{proof}

\section{Closing the integral}\label{sec:closing}

\begin{proof}[Proof of Theorem~\ref{thm:F}]
Split the integral into \(G_{\rm CL}\cup B_{\rm CL}\). On \(G_{\rm CL}\),
Lemma~\ref{lem:4.1}, \eqref{eq:int-DHDI}, and the finiteness
$\mu([\rstar,1])<\infty$ give
\[
   \int_{G_{\rm CL}\cap[\rstar,\rho_\delta]}
      |\bar z_\rho'|^2\,d\mu(\rho)
   \le C_*\int_{\rstar}^{1}(D_H+D_I)\,d\mu
       + C_*\,\delta_T(\Om)\,\mu([\rstar,1])
   \le C\delta.
\]
On \(B_{\rm CL}\), Lemma~\ref{lem:Lip} gives
\(|\bar z_\rho'|\le L_{\rm cent}\) a.e., while
\(\mu(B_{\rm CL})\le C\delta\). Hence
\[
   \int_{B_{\rm CL}\cap[\rstar,\rho_\delta]}
      |\bar z_\rho'|^2\,d\mu(\rho)
   \le L_{\rm cent}^2\mu(B_{\rm CL})\le C\delta.
\]
Combining the two estimates proves \eqref{eq:F}.\qedhere
\end{proof}

\section{Conditionality and remarks}\label{sec:conditionality}

\textbf{Claim: Theorem~\ref{thm:F} is unconditional.} It uses none of:
the per-$\rho$ pointwise inputs (1.1), (1.2)=(R), (2.4)
of the audit-trail markdown note (open in the no-graph regime);
the pointwise quadratic (B$^2$); Agent 3 (7.1); any
selection, graph entry, BDV minimiser, Schauder, or Fuglede expansion.

\emph{Structural reason:} the centroid kinematic identity~\eqref{eq:3.2}
has no homothety term, so the bound goes through the unconditional pair
(Cauchy--Schwarz on $V_\rho$; FMP+FJ+divergence at the centroid).

\begin{remark}[Sharpness]
The pointwise bound~\eqref{eq:4.1} is sharp in scaling: for a small
translation of a ball, $V_\rho-\bar V_\rho$ has dipole structure of
size $|\bar z_\rho'|$, giving $D_H\sim|\bar z_\rho'|^2$.
\end{remark}

\begin{remark}[Outside the smallness regime]
If $\delta_T(\Om)>\delta_0$, then $\mathcal A(\Om)\ge c(n)>0$ by FMP,
and \eqref{eq:F} is trivial after enlarging $C_{\rm cent}$.
\end{remark}

\begin{thebibliography}{99}

\bibitem{AFM2013}
E.~Acerbi, N.~Fusco, M.~Morini, CMP \textbf{322} (2013), 515--557.

\bibitem{AGS2008}
L.~Ambrosio, N.~Gigli, G.~Savar\'e, Birkh\"auser, 2008.

\bibitem{CL2012}
M.~Cicalese, G.\,P.~Leonardi, ARMA \textbf{206} (2012), 617--643.

\bibitem{FigalliMaggi2011}
A.~Figalli, F.~Maggi, ARMA \textbf{201} (2011), 143--207.

\bibitem{FMP2010}
A.~Figalli, F.~Maggi, A.~Pratelli, Invent. Math. \textbf{182} (2010), 167--211.

\bibitem{FJ2014}
N.~Fusco, V.~Julin, Calc. Var. PDE \textbf{50} (2014), 925--937.

\bibitem{Maggi2012}
F.~Maggi, Cambridge University Press, 2012.

\bibitem{Talenti1976}
G.~Talenti, Ann. Mat. Pura Appl. \textbf{110} (1976), 353--372.

\bibitem{LevelSetIdentity}
Building Block \textsc{LevelSetIdentity}, Plan~2.

\bibitem{MetricFramework}
Building Block \textsc{MetricFramework}, Plan~2.

\bibitem{SpaceTimeIdentity}
Building Block \textsc{SpaceTimeIdentity}, Plan~2.

\end{thebibliography}

\end{document}
